\newpage
\subsection{Constant eDOS Approximation}
\label{sec:constant_edos_approximation}

Previous section established that Eq.~\eqref{eq:intraband_constant_edos} can be evaluated numerically to high accuracy on a practical $\mathcal{E}$ grid. We now address the dominant approximation in the analytic pipeline the constant eDOS approximation (CEA), namely
\begin{equation}
\rho(\mathcal{E}+\hbar\omega)\rho(\mathcal{E})\;\approx\;\rho^2(\mathcal{E}_F),
\label{eq:constant_edos_replacement}
\end{equation}
which turns Eq.~\eqref{eq:intraband_energy_single} into Eq.~\eqref{eq:intraband_constant_edos}. The equations that we're going to compare here are the two that follow
\begin{equation}
\boxed{\begin{aligned}
\;& \int_{0}^{\infty} f^{T / S}(\mathcal{E}+\hbar\omega)\big[1-f^{T/S}(\mathcal{E})\big]\,
\rho(\mathcal{E}+\hbar\omega)\, \rho(\mathcal{E})\, d\mathcal{E} \\
\; &\int_{0}^{\infty}
 f^{T/S}(\mathcal{E}+\hbar\omega)\big[1-f^{T/S}(\mathcal{E})\big] \,\rho^{2}(\mathcal{E}_{F}) \;  d\mathcal{E}
\end{aligned}}
\label{eq:constant_edos_comparison_pair}
\end{equation}

The relative error will be calculated using the usual
\begin{equation}
\delta_{\mathrm{rel}}^{(3)}(\hbar\omega)
=\left|\frac{I_{(5)}(\hbar\omega)-I_{(4)}(\hbar\omega)}{I_{(4)}(\hbar\omega)}\right|
\label{eq:constant_edos_relative_error}
\end{equation}
where $I_{(4)}$ and $I_{(5)}$ denote the $\mathcal{E}$ integrals in Eq.~\eqref{eq:intraband_energy_single} and Eq.~\eqref{eq:intraband_constant_edos}, evaluated on the same numerical grid, and the T/S represent thermal and steady-state cases.

\subsubsection{Direct Relative Error of the Constant eDOS Approximation}
To separate the eDOS effect from the \emph{thermal} weighting, we also examine the point-wise mismatch
\begin{equation}
\delta_{\text{DOS}}(\mathcal{E},\hbar\omega)
=\left|\frac{\rho(\mathcal{E}+\hbar\omega)\rho(\mathcal{E})-\rho^{2}(\mathcal{E}_{F})}{\rho(\mathcal{E}+\hbar\omega)\rho(\mathcal{E})}\right|
=\left|1-\frac{\rho^{2}(\mathcal{E}_F)}{\rho(\mathcal{E}+\hbar\omega)\rho(\mathcal{E})}\right|.
\label{eq:constant_edos_pointwise}
\end{equation}

For the free-electron DOS $\rho(\mathcal{E})\propto\sqrt{\mathcal{E}}$ this reduces to
\begin{equation}
\delta_{\text{DOS}}(\mathcal{E},\hbar\omega)
=\left|1-\frac{\mathcal{E}_F}{\sqrt{\mathcal{E}(\mathcal{E}+\hbar\omega)}}\right|,
\label{eq:constant_edos_pointwise_reduced}
\end{equation}
so the locus of minimal DOS mismatch satisfies
\begin{equation}
\mathcal{E}(\mathcal{E}+\hbar\omega)=\mathcal{E}_F^2.
\label{eq:constant_edos_locus}
\end{equation}
Which can be seen in the figure below

\begin{figure}[H]
  \centering
  \includegraphics[width=0.85\textwidth]{edos_relative_error_default.png}
  \caption{Figure 3.1: Heat map of $\log_{10}\delta_{\text{DOS}}(\mathcal{E},\hbar\omega)$ for $\rho(\mathcal{E})\rho(\mathcal{E}+\hbar\omega)$ relative to $\rho^2(\mathcal{E}_F)$, evaluated for $\mathcal{E}_F=5\,\mathrm{eV}$.}
  \label{fig:edos_relative_error_default}
\end{figure}

Figure 3.1 shows that the constant-DOS-product assumption is only plausible in regions where the geometric mean $\sqrt{\mathcal{E}(\mathcal{E}+\hbar\omega)}$ remains close to $\mathcal{E}_F$. For fixed $\mathcal{E}_F$, large $\hbar\omega$ forces one of the two energies to be far from the other, and for small $\mathcal{E}$ the free-electron DOS varies rapidly and the mismatch necessarily grows.

\subsubsection{Emission Relative Error Due to Constant eDOS Approximation}

The physically relevant quantity, however, is the $\mathcal{E}$-integrated error (Figures 3.2 and 3.3), where $\delta_{\text{DOS}}(\mathcal{E},\hbar\omega)$ is weighted by the thermal factor $f^{T}(\mathcal{E}+\hbar\omega)[1-f^{T}(\mathcal{E})]$. Since this thermal factor localizes the integral to an energy window of width $\mathcal{O}(k_B T)$, the approximation quality depends on whether that window samples energies for which $\rho(\mathcal{E})$ is approximately constant on the relevant scale.

\begin{figure}[H]
  \centering
  \includegraphics[width=0.85\textwidth]{constant_eDOS_approximation.png}
  \caption{Figure 3.2 - Relative error in log-scale over the emission spectrum for several temperatures (thermal on the top steady-state in the bottom)}
  \label{fig:constant_edos_approximation}
\end{figure}

\begin{figure}[H]
  \centering
  \includegraphics[width=0.85\textwidth]{consteDOS_error_heatmaps.png}
  \caption{Figure 3.3 - Relative error in log-scale wrt the emission spectrum and temperature for both thermal and steady-state cases for a more comprehensive comparison}
  \label{fig:constedos_error_heatmaps}
\end{figure}

It appears as though the constant eDOS approximation is not very accurate. Need to discuss with Yonatan whether this is something he accepts, something I missed or perhaps something that should be fixed.

Having quantified the constant-eDOS approximation, the remaining stages address the earlier steps in the derivation---most notably the delta-function resolution that connects Eq.~\eqref{eq:intraband_energy_double} to Eq.~\eqref{eq:intraband_energy_single}.
